\RequirePackage{fix-cm}
\documentclass[smallextended,envcountsame]{svjour3}
\smartqed
\usepackage{amsmath,amssymb}
\usepackage{stix}
\usepackage{graphicx,booktabs,array,float,xcolor}
\usepackage[colorlinks=true,linkcolor=blue,citecolor=blue,urlcolor=blue]{hyperref}
\setlength{\emergencystretch}{2.5em}
\setlength{\tabcolsep}{4pt}
\newfloat{algorithm}{tbp}{loa}
\floatname{algorithm}{Algorithm}
\newcommand{\todo}[1]{\textcolor{red}{\textbf{[TODO: #1]}}}
\newcommand{\step}[1]{\par\noindent\hangindent=1.5em\hangafter=1\textbf{Step #1.}\ }
\newcommand{\repo}{\url{https://github.com/mossaiby/iHaar}}
\journalname{Journal of Scientific Computing}

\begin{document}

\title{Real spectra and fast tensor solvers for midpoint quadratic-spline collocation on graded grids}
\titlerunning{Real spectra and fast tensor solvers for quadratic-spline collocation}

\author{Farshid Mossaiby}
\authorrunning{F. Mossaiby}

\institute{F. Mossaiby \at
  Department of Civil Engineering, University of Isfahan, Isfahan 81744-73441, Iran\\
  \email{mossaiby@eng.ui.ac.ir}}

\date{Received: date / Accepted: date}

\maketitle

\begin{abstract}
Fast diagonalization solves separable problems on tensor grids with a few dense matrix products, but it
requires one-dimensional operators that are diagonalizable with real eigenvalues. Collocation operators
on graded grids are not symmetric, so this property has to be proved. We prove it for midpoint
quadratic-spline collocation, which is the nodal form of the Haar wavelet collocation method (HWCM). An
exact three-term identity writes the second-derivative operator as $A^{-1}L$ with tridiagonal $A$ and $L$.
Using the theory of oscillation matrices, we show that this operator has real, simple and nonpositive
eigenvalues on every strictly increasing grid, for any combination of prescribed values and zero slopes
at the ends. Fast diagonalization is therefore well defined on arbitrary graded tensor grids. We
distinguish this algebraic guarantee from backward stability, which we observe on random and adversarial
grids, and from forward accuracy, which is limited by the conditioning of the discrete problem. On
uniform grids, a correction that keeps the eigenvectors raises the order to four for zero-flux
conditions. Applications include mass-conserving Allen--Cahn and Cahn--Hilliard schemes with a modified
energy law, resolving interfaces with about four nodes across the transition layer, and adaptive plane
elastostatics for compressible materials. In a benchmark against biquadratic finite elements on
identical grids, the method needed about one quarter of the unknowns.
\keywords{Spline collocation \and Haar wavelet collocation \and Fast diagonalization \and Oscillation matrices \and Phase-field models \and Adaptive refinement}
\subclass{65N35 \and 65F15 \and 15B48 \and 65N22 \and 65M70 \and 74S25}
\end{abstract}

%==========================================================================
\section{Introduction}\label{sec:intro}
%==========================================================================
The Haar wavelet collocation method (HWCM)~\cite{chen1997haar,lepik2005numerical,lepik2011solving,lepik2014book}
is widely used in applied mechanics for beams, plates, shells and functionally graded
structures~\cite{majak2015accuracy,majak2015convergence,majak2018higher}. The highest derivative of the
unknown is expanded in Haar functions, the lower derivatives follow by exact integration, and the
equation is enforced at cell midpoints. Higher-order variants~\cite{majak2018higher,ratas2021solving} and
graded grids based on non-uniform Haar functions~\cite{dubeau2004nonuniform,lepik2008nonuniform,oruc2019nonuniform}
are available. In two dimensions, however, all equations are usually assembled into one dense matrix of
size $N^2\times N^2$, which becomes impractical beyond about $64\times64$ cells.

In terms of nodal values, HWCM is midpoint quadratic-spline collocation
(QSC)~\cite{houstis1988quadratic,christara1994quadratic,christara2006optimal}, and it uses the same
functions as quadratic isogeometric collocation~\cite{auricchio2010isogeometric,schillinger2013isogeometric}.
On tensor grids, separable problems can be solved by fast diagonalization~\cite{lynch1964direct}, which
reduces two- and three-dimensional problems to dense matrix products and fast transforms. It requires
each one-dimensional operator to be diagonalizable with real eigenvalues. For symmetric finite-element
and isogeometric Galerkin operators this is automatic~\cite{sangalli2016isogeometric}. Collocation
operators are generally not symmetric, and for them the property has to be proved. For Hermite-cubic
orthogonal spline collocation at Gauss points it is known~\cite{dyksen1987tensor,sun1999spectral}, and it
underlies the matrix decomposition algorithms of Bialecki and Fairweather~\cite{bialecki1995matrix,bialecki2011matrix}.
For QSC, fast solvers exist on uniform grids~\cite{christara2002fast}. To our knowledge, a spectral
guarantee has not been established for midpoint QSC on arbitrary non-uniform partitions. There, the
collocation operator is not symmetric and is not symmetrized by quadrature weights, so a different
argument is needed.

The main result of this paper (Theorem~\ref{thm:spectrum}) provides this guarantee. On every strictly
increasing grid, and for every combination of prescribed values and zero slopes at the ends, the
second-derivative operator has real, simple and nonpositive eigenvalues. The proof combines an exact
three-term identity (Theorem~\ref{thm:identity}) with the theory of oscillation
matrices~\cite{gantmacher2002oscillation,fallat2011totally}, and it is given in full in
Appendix~\ref{app:proofs}. We then examine what the result does and does not provide in floating-point
arithmetic, distinguishing well-definedness, backward stability and forward accuracy
(Section~\ref{sec:robust}). We also demonstrate its consequences in three settings: verification problems on
uniform and graded grids, Allen--Cahn and Cahn--Hilliard computations with thin interfaces, and
elastostatics with traction boundaries and adaptive refinement at corner singularities.
Table~\ref{tab:closest} compares the method with the closest existing approaches. Secondary proofs and
additional tables are collected in the Supplementary Material, and the code is available at \repo.
